% A name \let to a *character* is that character, everywhere a character is
% what the scanner wants. `\let\bgroup={` is how a document opens a group from
% inside a macro, and expl3's `\exp_stop_f:` is a space bound the same way --
% which is how it terminates a \romannumeral scan.
%
% The meaning was being stored and never dispatched, so every such name behaved
% as undefined the moment it was used.
\input plain
\catcode`\@=11
\def\M#1{\message{<#1>}}
\let\myb={ \let\mye=}
\setbox0=\hbox\myb aa\mye  \setbox1=\hbox{aa}
\M{brace=\the\wd0|\the\wd1}
% \let to a space needs two after the `=`: \let skips one optional space
% (tex.web 1221), so the second is what gets bound.
\def\usenn#1#2{#1#2}
\def\setsp#1{\let#1= }
\usenn{\setsp\spc}{ }
\M{meaning=\meaning\spc}
\setbox2=\hbox{a\spc b}\setbox3=\hbox{a b}
\M{space=\the\wd2|\the\wd3}
% And that is what stops a \romannumeral: the space ends the number scan.
\chardef\zed=0
\M{stopscan=\romannumeral\zed\spc ABC}
% \let to other character kinds.
\let\dollar=$ \let\amp=&
\M{other=\meaning\dollar|\meaning\amp}
% A \let-to-space terminates a *number* exactly as a space does. expl3 ends
% every \if_predicate:w and \if_int_compare:w in the kernel with \exp_stop_f:,
% which is bound this way, so a scanner that sees only a control sequence there
% cannot read any of them.
\count5=7
\M{ifnum=\ifnum\count5=7\spc yes\else no\fi}
\M{ifodd=\ifodd 3\spc yes\else no\fi}
\M{number=\number\count5\spc}
\M{relation=\ifnum 1<2\spc T\else F\fi}
\count0=5\spc \M{assign=\the\count0}
\M{numexpr=\the\numexpr 2+3\spc}

% §476 scans the resolved command code in a macro parameter text. A control
% sequence whose current meaning is a parameter character therefore introduces
% an argument; looking only at the token's written form gives this macro arity
% zero and makes the `#1` in its replacement text illegal.
\let\hash=#
\def\withimplicit[\hash1]{\M{macro-param=#1}}
\withimplicit[X]

% Braces use token identity, not resolved meaning, in a macro's parameter
% text. A control sequence let to either brace is therefore a legal delimiter;
% only an actual brace token starts/recovers the replacement scan. This is the
% distinction used by spectralsequences' smuggling macros.
\let\implicitopen={ \let\implicitclose=}
\def\untilclose#1\implicitclose{\M{close-delimiter=#1}}
\untilclose C\implicitclose
\def\untilopen#1\implicitopen{\M{open-delimiter=#1}}
\untilopen O\implicitopen
\def\implicitboth\hash1\implicitopen{\M{implicit-both=#1}}
\implicitboth B\implicitopen

% The `#{` exception still fires when the brace token really is `{`, including
% when the parameter character was reached through a control sequence.
\def\hashbrace\hash1\hash{\M{hash-brace=#1}}
\hashbrace H#{

% §§782--784 likewise inspect the resolved command while retaining the token
% itself in an alignment template. The glue scan expands `\peek` far enough to
% put its `\let` assignment into the u-template. At scan time `\probe` means
% `{`, but neither brace depth nor stored token identity may change: replay
% must still see the control sequence as the assignment target.
\let\probe={
\def\peek{\let\probe=*\M{align-token=preserved}}
\setbox4=\vbox{\halign{&\tabskip=-1pt\peek#\cr 1\cr}}
\bye
